\chapter{Технологическая часть}
%<писать все про реализацию: какие языки, какие ide и тд.; коды алгоритмов/программы>
%<тестовые данные и волшебные слова <<все тесты пройдены успешно>> >

\section{Реализация алгоритмов}
Для реализации алгоритмов использовался язык программирования C++, так как он поддерживает статическую типизацию и имеет возможность отключения сбощика мусора. Отладка происходила в среде разработки <<Visual Studio 2022>>.\\

% 14911

\lstinputlisting[firstline=10, lastline=17, style=cpp]{../code/recursive.cpp}
Листинг 1 -- Рекурсивный алгоритм.\\
{}

\lstinputlisting[firstline=5, lastline=8, style=cpp]{../code/recursive.cpp}
Листинг 2 -- Нерекурсивный алгоритм.\\
{}


\section{Тестирование}

\begin{table}[!htbp]
	\caption{Модульные тесты для рекурсивного алгоритма.}
	\centering
	\begin{tabularx}{\textwidth}{|*{2}{>{\centering\arraybackslash}X|}}
		\hline
		\textbf{Вход} & \textbf{Выход}\\\hline
		1 & 1 \\\hline
		10 & 1 2 3 4 5 6 7 8 9 10 \\\hline
	\end{tabularx}
\end{table}

\begin{table}[!htbp]
	\caption{Модульные тесты для нерекурсивного алгоритма.}
	\centering
	\begin{tabularx}{\textwidth}{|*{2}{>{\centering\arraybackslash}X|}}
		\hline
		\textbf{Вход} & \textbf{Выход}\\\hline
		1 & 1 \\\hline
		10 & 1 2 3 4 5 6 7 8 9 10 \\\hline
	\end{tabularx}
\end{table}


\section*{Вывод}
%<что делали и получили в результате>

В данном разделе мной были реализованы рекурсивная и нерекурсивная версии алгоритмов по варианту и проведено их успешное тестирование.

\clearpage
